\section{Atomic Physics}

\begin{problem}{5.01}
A clean metal surface is irradiated with ultraviolet light of wavelength $\lambda = 200\unitv{nm}$, and the stopping voltage of the photoelectrons is measured to be $2.6\unitv{V}$. What is the work function of this metal? If ultraviolet light of wavelength $\lambda = 300\unitv{nm}$ is used instead, what is the stopping voltage? What happens if visible light is used? ($hc = 1240\unitv{eV\cdot nm}$)
\end{problem}

\begin{problem}{5.02}
In order to decide whether a particle of very great penetrating power is a $\gamma$ photon, scientists carried out the following work: (a) they measured the attenuation of this ``photon'' in lead and estimated its energy to be about $15\unitv{MeV}$; (b) they bombarded paraffin ($\mathrm{CH_4}$) with a beam of these particles, producing a large number of protons, the maximum recoil energy of the protons being $5.7\unitv{MeV}$. From these two experimental results, decide by calculation whether or not this particle is a $\gamma$ photon. (Electron rest energy $m_ec^2 = 0.511\unitv{MeV}$, $hc = 12\,400\unitv{eV}\cdot\text{\AA}$)
\end{problem}

\begin{problem}{5.03}
(1) Find the Compton shift $\Delta\lambda$ in the following two cases, the photon being incident on a free electron initially at rest:
(a) if after scattering the electron moves along the direction of incidence of the photon;
(b) if the photon still moves along its original direction of incidence.

(2) A photon collides with a high-energy relativistic electron; the photon is scattered through an angle of $60\degree$, while the electron comes to rest.
(a) Find the wavelength difference $\Delta\lambda = \lambda - \lambda'$ between the incident photon (wavelength $\lambda$) and the outgoing photon (wavelength $\lambda'$).
(b) If the energy of the incident photon is equal to the rest energy of the electron, find the energy of the scattered photon and the kinetic energy of the electron before the collision.
\end{problem}

\begin{problem}{5.04}
SLAC (Stanford Linear Accelerator Center) once used the following scheme to produce high-energy monochromatic photons: a low-energy photon beam from a laser (in the visible region) collides head-on with a $20\unitv{GeV}$ electron beam.
\begin{enumerate}[label=(\arabic*),leftmargin=2.2em,itemsep=1pt,topsep=2pt]
\item Let $\theta$ be the angle between the scattered photon beam and the incident electron beam. Find the energy of the scattered photon as a function of $\theta$.
\item Give an approximate value for the maximum energy of the scattered photon (in GeV).
\end{enumerate}
\end{problem}

\begin{problem}{5.05}
The track of a $1\unitv{MeV}$ electron in a cloud chamber is a string of small water droplets, the linear size of a droplet being about $10^{-5}\unitv{m}$. Will the track of the electron differ appreciably from a straight line? Why? (Electron rest energy $m_ec^2 = 0.511\unitv{MeV}$, $hc = 12\,400\unitv{eV}\cdot\text{\AA}$)
\end{problem}

\begin{problem}{5.06}
An ion has nuclear charge number $z$ and only one electron of rest mass $m$ outside the nucleus. Using the uncertainty relation $\Delta x\,\Delta p \gtrsim \hslash$, estimate the ground-state energy of this hydrogen-like ion in the following two cases (1) and (2).

(1) In the non-relativistic case, when the electron is at a distance $r$ from the nucleus and has momentum $p$, the energy of the ion is
\[
E = \frac{p^2}{2m} - \frac{ze^2}{4\pi\varepsilon_0 r}.
\]

(2) In the relativistic case, the energy of the ion is
\[
E = \sqrt{p^2c^2 + m^2c^4} - \frac{ze^2}{4\pi\varepsilon_0 r}.
\]

(3) What is the connection between the results of (1) and (2)?
\end{problem}

\begin{problem}{5.07}
Calculate:
\begin{enumerate}[label=(\arabic*),leftmargin=2.2em,itemsep=1pt,topsep=2pt]
\item the ratio of the circumference of the $n$-th orbit of the hydrogen atom to the de Broglie wavelength of the electron moving in that orbit;
\item the de Broglie wavelength of the electron moving in the $n=1$ orbit of the hydrogen atom.
\end{enumerate}
\end{problem}

\begin{problem}{5.08}
Because the hydrogen gas contained helium gas, the spectral lines of the helium ion ($\mathrm{He}^+$) were mixed into the photographed spectrum of atomic hydrogen. Show how, without measuring the wavelengths of the spectral lines, one can decide directly from the photograph which lines belong to the hydrogen spectrum and which to the $\mathrm{He}^+$ spectrum, and briefly explain the basis of your judgement.
\end{problem}

\begin{problem}{5.09}
The $\mu^-$ meson carries a charge $-e$ and has mass $m_\mu = 106\unitv{MeV}/c^2$; the $\pi^+$ meson carries a charge $+e$ and has mass $m_\pi = 140\unitv{MeV}/c^2$. If in the ``exotic atom'' formed by these two mesons both particles may be regarded as point particles, find:
\begin{enumerate}[label=(\arabic*),leftmargin=2.2em,itemsep=1pt,topsep=2pt]
\item the lowest energy of the atom;
\item the distance between the two mesons in the ground state;
\item the speed of the $\mu^-$ meson in the ground state, neglecting the overall motion of the atom.
\end{enumerate}
\end{problem}

\begin{problem}{5.10}
A $\mu^-$ meson is captured by a lead nucleus ($^{208}_{82}\mathrm{Pb}$, $z = 82$), forming a lead muonic atom. Apart from its different mass, the $\mu^-$ meson is the same thing as an electron; the mass of the $\mu^-$ meson is known to be $207$ times the electron mass. Assume that the nuclear charge is uniformly distributed within a sphere of radius $R = 7.7\unitv{fm}$ and that the $\mu^-$ meson can move inside this sphere. Using Bohr's theory, find:
\begin{enumerate}[label=(\arabic*),leftmargin=2.2em,itemsep=1pt,topsep=2pt]
\item the radii of the first two Bohr orbits of the $\mu^-$ meson;
\item the energies of the first two levels of the lead muonic atom;
\item the wavelength of the first line of the Lyman series of the lead muonic atom;
\item how the result of part (3) changes if the lead nucleus is treated as a point charge.
\end{enumerate}
\end{problem}

\begin{problem}{5.11}
The sodium atom ($z = 11$) has $11$ electrons outside the nucleus, of which $10$ are rather tightly bound to the nucleus and form the atomic core, while the outermost valence electron is rather loosely bound to the core. In the spectrum emitted when the valence electron of the sodium atom makes a transition, a spectral series similar to the Lyman series of hydrogen is described by the formula
\[
\frac{1}{\lambda} = 0.378R - \frac{R}{m^2},
\]
where $R$ is the Rydberg constant and $m$ is an integer.
\begin{enumerate}[label=(\arabic*),leftmargin=2.2em,itemsep=1pt,topsep=2pt]
\item The formula holds only within a certain range; explain, in terms of atomic structure, the conditions under which the formula holds.
\item Use this formula to calculate the ionisation energy of sodium (that is, the energy required to ionise the valence electron).
\end{enumerate}
\end{problem}

\begin{problem}{5.12}
Assume that Bohr's theory of the hydrogen atom applies to each of the following parts, and neglect the motion of the helium nucleus.
\begin{enumerate}[label=(\arabic*),leftmargin=2.2em,itemsep=1pt,topsep=2pt]
\item Find the ground-state energy of the helium ion ($\mathrm{He}^+$) and the radius of the first Bohr orbit of the electron.
\item The mass of the $\mu^-$ meson is about $200$ times the electron mass; apart from the mass, it is identical with the electron in every respect. Suppose a helium nucleus captures a $\mu^-$ meson and forms a hydrogen-like ion. Find the ground-state energy of this system and the radius of the first Bohr orbit of the $\mu^-$ meson.
\item Suppose the atom formed when one electron of the helium atom is replaced by a $\mu^-$ meson is in its ground state. (a) Estimate the radius of the first Bohr orbit of the electron; (b) estimate the energy of this system; (c) to which chemical element is the chemical behaviour of this atom most similar, and why?
\end{enumerate}
(Numerical results in this problem need only two significant figures; energies are to be given in $\unitv{eV}$ and radii in $\unitv{nm}$.)
\end{problem}

\begin{problem}{5.13}
Hydrogen atoms are bombarded with particles and thereby excited. In each of the following cases, what is the minimum kinetic energy that the particle must have in order that the hydrogen atom emit the $\mathrm{H}_\alpha$ line?
\begin{enumerate}[label=(\arabic*),leftmargin=2.2em,itemsep=1pt,topsep=2pt]
\item An electron collides with a hydrogen atom at room temperature.
\item A proton bombards a ground-state hydrogen atom initially at rest.
\item A proton and a ground-state hydrogen atom approach each other with equal speeds and collide head-on.
\end{enumerate}
\end{problem}

\begin{problem}{5.14}
The helium atom is the simplest many-electron atom.
\begin{enumerate}[label=(\arabic*),leftmargin=2.2em,itemsep=1pt,topsep=2pt]
\item If there were no interaction between the electrons, what would be the ground-state energy of the helium atom?
\item Using a rough model, estimate the ground-state energy of the helium atom when the Coulomb interaction between the electrons is taken into account.
\item If one electron of the helium atom is excited to the $4f$ state while the other electron remains in the $1s$ state, what is the energy of the helium atom then?
\end{enumerate}
\end{problem}

\begin{problem}{5.15}
The natural radioactive uranium series starts from $^{238}_{92}\mathrm{U}$, the uranium isotope with the longest half-life, and after a series of $\alpha$ and $\beta$ decays finally reaches the stable nucleus $^{206}_{82}\mathrm{Pb}$. The relative ratio of the numbers of $^{238}_{92}\mathrm{U}$ and $^{206}_{82}\mathrm{Pb}$ nuclei measured at present in uranium ore is
\[
\frac{N(^{238}_{92}\mathrm{U})}{N(^{206}_{82}\mathrm{Pb})} = 2.785.
\]
Estimate the age of the Earth. (Given: the half-life of $^{238}_{92}\mathrm{U}$ is $T_{1/2} = 4.468\times10^9$ years, and among its daughter nuclei $^{234}_{92}\mathrm{U}$ has the longest half-life, $T_{1/2} = 2.45\times10^5$ years.)
\end{problem}

\begin{problem}{5.16}
It is known that the concentration ratio $^{14}\mathrm{C}/^{12}\mathrm{C}$ in the atmosphere is $\eta = 1.2\times10^{-12}$, and that the half-life of $^{14}\mathrm{C}$ is $T_{1/2} = 5730$ years. Since a living organism must necessarily exchange carbon with the atmosphere through respiration or photosynthesis, the content ratio of $^{14}\mathrm{C}$ to $^{12}\mathrm{C}$ in a living body is the same as in the atmosphere.
\begin{enumerate}[label=(\arabic*),leftmargin=2.2em,itemsep=1pt,topsep=2pt]
\item Calculate the specific activity of the carbon in the atmosphere, that is, the number of disintegrations per minute of one gram of carbon.
\item The $^{14}\mathrm{C}$ content of an ancient biological sample is $1/1.831$ of that of a living biological sample. Estimate the age of this ancient sample.
\item The relative statistical error of a radioactive count $N$ is defined as $\delta = 1/\sqrt{N}$. In order that the statistical error of the measurement be $\delta = 1\%$, how many grams of the above ancient sample are required so that the count over one hour satisfies this requirement on $\delta$? Also estimate the precision, plus or minus how many years, of the age determined at this $\delta$. (Assume that the efficiency of the detector is $100\%$.)
\end{enumerate}
\end{problem}

\begin{problem}{5.17}
(1) Write down the expression for $\beta^-$ decay and the formula for its decay energy (the decay-energy formula is required to be expressed in terms of atomic masses).

(2) As the figure shows, $^{17}\mathrm{N}$ decays by $\beta^-$ emission, with a maximum electron kinetic energy of $3.72\unitv{MeV}$, to an excited state of $^{17}\mathrm{O}$ ($^{17}\mathrm{O}^*$), and from this excited state it decays further by emitting a neutron to the ground state or to an excited state of $^{16}\mathrm{O}$. It is known that $^{16}\mathrm{O}$ has several excited states ($^{16}\mathrm{O}^*$) at $3.05$, $6.13$, $6.90\unitv{MeV}$, etc. (the ground-state energy of $^{16}\mathrm{O}$ is taken as $0$). State how many possible values the kinetic energy of the emitted neutron can take, and find the corresponding values. The recoil of the $^{16}\mathrm{O}$ nucleus when the neutron is emitted must be taken into account in the calculation.

(Known: atomic masses $M(^{17}\mathrm{N}) = 17.008\,449\unitv{u}$, $M(^{16}\mathrm{O}) = 15.994\,915\unitv{u}$, neutron mass $m_n = 1.008\,665\unitv{u}$, $1\unitv{u}\,c^2 = 931.5\unitv{MeV}$, where $\unitv{u}$ is the atomic mass unit.)
\end{problem}

\begin{figbox}{5.17}
\begin{tikzpicture}[scale=1.0]
  % energy levels
  \draw[thickline] (1.2,4.6) -- (4.2,4.6) node[right] {$^{17}\mathrm{N}$};
  \draw[thickline] (2.3,3.35) -- (6.0,3.35) node[right] {$^{17}\mathrm{O}^*$};
  \draw[thickline] (5.2,2.25) -- (7.0,2.25) node[right] {$^{16}\mathrm{O}^*$};
  \draw[thickline] (3.8,1.15) -- (7.2,1.15) node[right] {$^{16}\mathrm{O}$};
  % transitions
  \draw[->,line width=0.7pt] (1.9,4.53) -- (2.78,3.40);
  \node at (2.66,4.06) {$\beta^-$};
  \draw[->,line width=0.7pt] (4.95,3.30) -- (5.55,2.30);
  \node at (5.70,2.82) {$n'$};
  \draw[->,line width=0.7pt] (2.35,3.30) -- (4.05,1.20);
  \node at (2.76,2.40) {$n$};
\end{tikzpicture}\\[2pt]
{\small In the figure $^{16}\mathrm{O}^*$ denotes an excited state; in reality it is not a single level. $n'$ denotes the neutron emitted in the decay to this excited state, and its kinetic energy differs from that of the neutron $n$ emitted in the decay to the ground state.}
\end{figbox}

\begin{problem}{5.18}
\begin{enumerate}[label=(\arabic*),leftmargin=2.2em,itemsep=1pt,topsep=2pt]
\item Write down the semi-empirical formula for the binding energy of a nucleus in the liquid-drop model, and state the physical meaning of each term and which terms come from quantum effects.
\item The atomic masses of a pair of mirror nuclei (that is, two nuclei with the same mass number $A$ but with the proton and neutron numbers interchanged) --- boron and carbon --- are known to be $M(^{11}\mathrm{B}) = 11.009\,31\unitv{u}$ and $M(^{11}\mathrm{C}) = 11.011\,43\unitv{u}$. Use the semi-empirical binding-energy formula to find the value of the constant $r_0$ in the corresponding nuclear-radius formula $R = r_0A^{1/3}$.
\end{enumerate}
(Known: hydrogen atom mass $M(^{1}\mathrm{H}) = 1.007\,825\unitv{u}$, neutron mass $m_n = 1.008\,665\unitv{u}$, combination constant $e^2/4\pi\varepsilon_0 = 1.44\unitv{MeV\cdot fm}$, $1\unitv{u}\,c^2 = 931.5\unitv{MeV}$.)
\end{problem}

\begin{problem}{5.19}
In the Rutherford elastic-scattering experiment, the scattering of the incident particle and the recoil of the target nucleus occur together (see the figure). In the figure, $\theta_L$ and $\varphi_L$ denote respectively the scattering angle and the recoil angle of the target nucleus in the laboratory system, and $\theta_C$ and $\varphi_C$ denote the corresponding angles in the centre-of-mass system. Assume that the incident particle is an alpha particle ($^{4}_{2}\mathrm{He}$) of kinetic energy $2\unitv{MeV}$ bombarding a hydrogen nucleus ($^{1}_{1}\mathrm{H}$) assumed to be at rest in the target, and carry out the following calculations:
\begin{enumerate}[label=(\arabic*),leftmargin=2.2em,itemsep=1pt,topsep=2pt]
\item What is the kinetic energy of the recoiling hydrogen nucleus in the direction $\varphi_L = 30\degree$ in the laboratory system?
\item What is the Rutherford differential scattering cross-section of the recoiling hydrogen nucleus at $\varphi_L = 30\degree$ in the laboratory system?
\end{enumerate}
Given: in the centre-of-mass system, the Rutherford differential scattering cross-section for the outgoing particle is
\[
\sigma_C(\theta_C) = \left(\frac{a}{4}\right)^2 \frac{1}{\sin^4(\theta_C/2)},
\]
where the Coulomb scattering factor is
\[
a = \frac{z_1z_2e^2}{4\pi\varepsilon_0E_C},
\]
$z_1$ and $z_2$ being respectively the charge numbers of the incident particle and of the target nucleus, and $E_C$ the centre-of-mass energy of the incident particle. The combination constant is $e^2/(4\pi\varepsilon_0) = 1.44\unitv{eV\cdot nm}$.
\end{problem}

\begin{figbox}{5.19}
\begin{tikzpicture}[scale=0.95]
  \coordinate (T) at (3.0,3.0);
  \coordinate (C) at (6.4,3.0);
  \coordinate (P) at (5.3,4.7);
  \coordinate (R) at (8.3,1.0);
  % beam axis
  \draw[thickline,->] (1.10,3.0) -- (2.74,3.0);
  \draw[dashedline] (3.24,3.0) -- (9.2,3.0);
  \draw[force] (5.95,3.0) -- (6.30,3.0);
  % outgoing tracks
  \draw[thickline,->] (3.13,3.10) -- (5.06,4.52);
  \draw[thickline,->] (3.15,2.94) -- (8.14,1.06);
  % centre-of-mass construction
  \draw[dashedline] (C) -- (5.39,4.565);
  \draw[dashedline] (C) -- (8.19,1.116);
  % particles
  \draw[thickline,fill=white] (0.90,3.0) circle (0.16);
  \draw[thickline,fill=white] (T) circle (0.16);
  \draw[thickline,fill=white] (P) circle (0.16);
  \draw[thickline,fill=white] (R) circle (0.16);
  % angles at the target
  \draw[thinline] ($(T)+(1.00,0)$) arc (0:36.5:1.00);
  \node at (4.24,3.42) {$\theta_L$};
  \draw[thinline] ($(T)+(0.85,0)$) arc (0:-20.7:0.85);
  \node at (4.28,2.70) {$\varphi_L$};
  % angles at the centre of mass
  \draw[thinline] ($(C)+(122.9:0.80)$) arc (122.9:180:0.80);
  \node at (5.44,3.54) {$\theta_C$};
  \draw[thinline] ($(C)+(-46.5:0.80)$) arc (-46.5:0:0.80);
  \node at (7.42,2.58) {$\varphi_C$};
  % labels
  \node[below] at (0.90,2.82) {\small $M(^{4}_{2}\mathrm{He})$};
  \node[below] at (2.80,2.72) {\small $M(^{1}_{1}\mathrm{H})$};
  \node[below left] at (6.38,2.98) {$C$};
\end{tikzpicture}\\[2pt]
{\small Elastic scattering of an alpha particle ($^{4}_{2}\mathrm{He}$) by a hydrogen nucleus; $C$ is the centre of mass of the two nuclei.}
\end{figbox}
